% C097: 简要导入材料 / source brief
\noindent\textbf{C097 抛物线内接三角形的斜率关系}\par
\noindent 设 $p\ne0$，$O$ 为坐标原点，$A,B$ 为抛物线上两个不同的非顶点。对两种标准形式，有如下对应关系：
\[
\boxed{y^2=2px\quad\Longrightarrow\quad
\frac1{k_{OA}}+\frac1{k_{OB}}=\frac1{k_{AB}}}\qquad
\boxed{x^2=2py\quad\Longrightarrow\quad
k_{OA}+k_{OB}=k_{AB}}.
\]
\noindent\textbf{边界：}两式按通常斜率理解时，相关斜率必须存在；对 $y^2=2px$，若 $AB$ 是竖直直线，则 $k_{AB}$ 不存在，不能写 $1/k_{AB}$。更通用的写法是
\[
\frac1{k_{OA}}+\frac1{k_{OB}}=\frac{x_B-x_A}{y_B-y_A},
\]
其中右边在 $A\ne B$ 时始终有定义。对 $x^2=2py$，$AB$ 不会竖直，且 $k_{AB}=0$ 是允许的。
\noindent\textbf{最快记忆：}横向开口时\emph{倒斜率相加}，纵向开口时\emph{斜率相加}；两者都是平方差公式 $(u^2-v^2)=(u-v)(u+v)$ 的直接结果。\par
\noindent\textbf{关联：}正文保留四种证明，并补充一般二次函数的顶点平移、两端切线的斜率平均关系、倒数和快速变形；焦点弦长度倒数和及直角三角形倒数平方和仅作结构对照。\par
